\documentclass[12pt]{article}

\usepackage[left=1.25in,right=1.25in,top=1in,bottom=1in]{geometry}
\usepackage{setspace}
\usepackage{microtype}
\usepackage{amsmath,amssymb,amsfonts}
\usepackage{mathtools}
\usepackage{graphicx}
\usepackage{hyperref}

\setstretch{1.15}
\setlength{\parindent}{0pt}
\setlength{\parskip}{0.75em}

\title{The Politics of Oracle-Free Evaluation:\\
Constraint, Generativity, and the Problem of Validation}
\author{Flyxion}
\date{\today}

\begin{document}

\maketitle

\begin{abstract}

Recent proposals for ``oracle-free'' evaluation seek to validate mathematical reasoning without access to ground-truth answers by measuring how well candidate solutions enable the resolution of related problems. Consequence-Based Utility (CBU) represents the most explicit formulation of this approach, replacing verification by correspondence with verification by generativity. Instead of asking whether a solution is correct, it asks whether the solution behaves as if it were correct when used as an in-context exemplar (Son et al., 2026).

This paper argues that such methods are reliable only within closed formal domains and become systematically misleading when extended to open, embodied, or temporally evolving systems. I introduce the notion of \emph{constraint-indexed correctness}, according to which the validity of an explanation is determined by its stability under transformations imposed by physical law, scale transition, measurement, and historical development. The analysis situates oracle-free evaluation within a longer shift from reconstructive validation toward operational metrics (Kuhn, 1962; Lakatos, 1976; Hacking, 1983) and examines how contemporary infrastructures redistribute epistemic authority (Lessig, 1999; Marres, 2017).

The goal is not to reject consequence-based reasoning but to reposition it as an inner-loop heuristic embedded within broader constraint-sensitive validation. Scientific knowledge advances not through unconstrained generativity, but through sustained friction with reality.

\end{abstract}

\section{Introduction}

The ambition to evaluate reasoning without recourse to authoritative answers has deep roots in both mathematics and computation. If solutions could be validated without explicit ground truth, inquiry might scale indefinitely into domains where verification is costly or delayed. Oracle-free evaluation embodies this aspiration by attempting to infer correctness indirectly, through observable consequences of adopting a candidate explanation.

Consequence-Based Utility provides a concrete instantiation. A proposed solution is not judged in isolation but by the extent to which it enables successful reasoning on nearby problems. Correctness is inferred from downstream usefulness rather than direct proof (Son et al., 2026).

At first glance, this appears to mirror ordinary mathematical practice. Mathematicians often gain confidence in arguments that unify disparate results or generate productive lines of inquiry. Yet historical studies emphasize that such confidence emerges through prolonged communal reconstruction, critique, and reinterpretation (Lakatos, 1976). Validation is not merely inferential but social, material, and temporal. Scientific authority rests upon reproducibility, shared norms, and durable records of reasoning (Merton, 1973; Shapin, 1994).

Oracle-free evaluation compresses this extended process into short computational loops. Instead of reconstruction by a community, validation is delegated to performance signals measured within a designed environment. The transition marks a shift from verification as constraint to verification as throughput.

The question, therefore, is not whether such systems function, but what kind of correctness they detect.

\section{From Correspondence to Operational Adequacy}

Classical philosophy of science distinguished between correspondence with reality and coherence within a theoretical structure. In mathematics these often coincide because the objects under study are themselves formal. In empirical domains the distinction becomes decisive.

Operational validation risks collapsing this difference. When correctness is inferred from predictive success or inferential utility alone, evaluation becomes vulnerable to what has been described as the failure of proxy metrics to track intended goals (Goodhart, 1984).

Cartwright's analysis of scientific laws showed that successful models operate only within constrained regimes rather than universally (Cartwright, 1983). Anderson's work on emergence demonstrated that explanations valid at one scale can fail catastrophically at another (Anderson, 1972). These insights imply that correctness must be indexed to the conditions under which reasoning remains stable.

Consequence-based validation measures coherence within a constructed manifold. Constraint-based validation tests invariance under transformations imposed by the world.

\section{The Structure of Consequence-Based Evaluation}

CBU replaces a binary oracle of correctness with a performance functional. Let \(Q\) denote a problem, \(C\) a candidate solution, and \(\mathcal{N}(Q)\) a curated set of related questions. A solver conditioned on \(C\) is evaluated over \(\mathcal{N}(Q)\), and its accuracy becomes a proxy for the validity of \(C\) (Son et al., 2026).

This procedure measures what may be called \emph{method-level transportability}. A solution is valued insofar as it encodes reusable reasoning patterns.

Yet the construction presupposes that correctness is equivalent to such transportability. This assumption is secure only in domains whose structure is closed under the transformations being sampled. Once explanations must survive scale change, measurement, or causal intervention, transferability becomes insufficient.

The remainder of this paper develops a framework for understanding these limitations and articulates an alternative grounded in constraint-indexed validation.

\section{Constraint-Indexed Correctness}

To evaluate explanations in open systems, correctness must be defined relative to the constraints under which those systems operate. I therefore introduce the notion of \emph{constraint-indexed correctness}, which replaces neighborhood propagation with invariance under admissible transformations.

Let
\[
\mathcal{K} = (\Phi, \Sigma, \Tau, \Mu)
\]
be a tuple of constraint classes. The component \(\Phi\) represents invariant relations, such as conservation laws or definitional identities. The component \(\Sigma\) denotes scale transformations mapping descriptions across levels of resolution. The component \(\Tau\) captures temporal evolution operators governing system dynamics. The component \(\Mu\) specifies measurement interfaces through which theoretical quantities become operationally accessible.

Given a candidate explanation \(C\) addressing a claim \(Q\), we say that \(C\) is correct relative to \(\mathcal{K}\) if its inferential structure remains stable under all admissible transformations generated by these constraints. Formally,
\[
C \models_{\mathcal{K}} Q
\]
when for every transformation \(T\) induced by \(\mathcal{K}\), the transformed explanation \(T(C)\) entails the appropriately transformed claim \(T(Q)\).

Where consequence-based evaluation samples nearby formal questions, constraint-indexed validation samples nearby \emph{realizations}. It asks not whether reasoning generalizes syntactically, but whether it survives interaction with scale, causality, and measurement.

This view resonates with interventionist accounts of explanation, in which understanding is tied to stability under manipulation rather than descriptive adequacy alone (Hacking, 1983; Pearl, 2009).

\section{Local Transfer Versus Structural Stability}

The difference between consequence-based and constraint-indexed validation can be illustrated by considering candidates that behave identically within a curated neighborhood yet diverge under constraint transformation.

Suppose two derivations \(C_1\) and \(C_2\) yield equivalent predictions across a set of nearby problems constructed by perturbation. A consequence-based evaluator cannot distinguish them. Yet if \(C_2\) violates an invariant encoded in \(\Phi\), repeated application of a scale transformation \(S \in \Sigma\) reveals systematic divergence.

Such behavior is common in scientific modeling. Linear approximations function locally but fail under renormalization or regime change. Mean-field theories capture average behavior yet collapse near critical transitions. These phenomena reflect the general principle that explanations must remain stable across scales rather than merely across analogous statements (Anderson, 1972).

Consequence-based validation detects coherence within a local manifold. Constraint-indexed validation detects compatibility with the structure that generates the manifold itself.

\section{Temporal Embedding}

CBU evaluates candidate reasoning through bounded rollouts, implicitly treating correctness as time-independent. Yet explanations in open systems must remain reliable as conditions evolve.

Let observational updates be indexed by time \(t\). Define a temporal validation functional
\[
U_t(C) = \mathbb{E}\bigl[\text{performance of } C \text{ under constraints observed up to } t\bigr].
\]
Correctness becomes trajectory-dependent rather than snapshot-based. Static consequence-based scoring appears as a limiting case obtained only when explanatory validity is independent of temporal evolution, an assumption rarely satisfied in practice.

Historical analyses of science emphasize precisely this temporal dimension: theories persist by surviving successive confrontations with new data rather than by performing well in isolated evaluations (Kuhn, 1962).

\section{Embodiment and Instantiation}

Formal reasoning incurs negligible marginal cost. Instantiated explanations do not. Real systems impose thermodynamic, material, and measurement constraints that eliminate vast regions of formally admissible structure.

Let \(\mathcal{I}\) denote an instantiation operator mapping abstract descriptions into realizable processes. A candidate explanation is viable only if \(\mathcal{I}(C)\) exists and preserves predictive coherence under physical constraints.

This emphasis aligns with accounts of knowledge as materially situated practice rather than purely symbolic manipulation (Polanyi, 1966; Sennett, 2008; Crawford, 2009). Constraint-indexed validation treats embodiment as a primary epistemic filter rather than an implementation detail.

\section{Information, Compression, and External Reference}

Consequence-based evaluation may also be interpreted through information theory. It measures how efficiently a candidate explanation compresses reasoning within a constructed neighborhood. Constraint-indexed validation instead asks whether this compression preserves mutual information with an external system.

Efficient encoding alone does not guarantee semantic fidelity (Cover and Thomas, 2006). A model may generalize internally while discarding structure required for intervention or extrapolation. Validation must therefore assess preservation of externally anchored information, not merely internal consistency.

\section{Infrastructure as Epistemology}

Evaluation frameworks are inseparable from the infrastructures that sustain them. Architectural choices shape what counts as evidence, who can participate, and how claims are reproduced. Digital systems regulate behavior through design in ways analogous to legal or institutional rules (Lessig, 1999).

Studies of platform-mediated knowledge production show how evaluation becomes entangled with metrics, visibility, and access (Marres, 2017). Earlier scientific regimes relied upon durable archives and independent reproducibility as conditions of credibility (Merton, 1973). When validation migrates into tightly coupled computational ecosystems, epistemic authority shifts accordingly.

The issue is not technological change but the relocation of judgment from reconstructive practice to infrastructural mediation.

\section{Adversarial Considerations}

A critique of consequence-based evaluation must address its strongest defense. Within strictly formal domains, one may argue that consequence is indistinguishable from truth. If mathematical objects are defined entirely by their inferential roles, then the ability of a candidate solution to organize correct reasoning across related cases captures precisely what correctness means.

This position has historical precedent in structuralist interpretations of mathematics, where objects are understood as positions in relational systems rather than as independently existing entities. Under such assumptions, validation through coherent extension is not a proxy but an expression of mathematical validity itself.

The present argument does not deny this equivalence. Instead, it asserts that the equivalence depends upon closure conditions rarely satisfied outside pure formalism. Once explanations interact with measurement, scale transition, or embodiment, inferential fertility ceases to guarantee adequacy.

A second objection holds that science itself often advances through coherence rather than direct falsification. Research programs adapt to anomalies rather than collapsing immediately (Lakatos, 1976). Yet such adaptation occurs within an environment of empirical resistance. Adjustments must eventually reconcile with observation or fail to stabilize historically.

A third objection notes that constraint-indexed validation also depends on human judgment in specifying constraints. This is correct. Constraints are neither self-evident nor immutable. The distinction lies in whether such assumptions remain implicit within datasets and evaluation protocols, or are explicitly modeled and subjected to revision.

\section{Computational Cost and the Scalability Question}

Consequence-based methods scale efficiently because they reduce validation to repeated inference within a fixed model. Constraint-indexed approaches impose additional demands, requiring analysis of invariance under transformation, simulation, or measurement.

This asymmetry reflects a broader historical pattern. Empirical validation has always been more resource-intensive than formal speculation. Instrumentation, replication, and cross-scale analysis impose costs absent from purely symbolic reasoning.

The relevant question is not whether constraint-aware validation is cheaper, but whether its additional expense corresponds to epistemic value. If correctness in open systems depends on confrontation with constraint, then eliminating that confrontation risks substituting efficiency for reliability.

\section{The Bootstrap Problem of Constraint Discovery}

Constraint-indexed validation presupposes knowledge of the constraint tuple \(\mathcal{K}\). Yet such knowledge is itself provisional. Scientific practice resolves this apparent circularity iteratively. Constraints are discovered, refined, and occasionally abandoned through cycles of prediction and failure.

Theories do not begin with perfect constraints; they acquire them through engagement with experiment and instrumentation. Validation is therefore not static but historically extended, an interaction between conceptual structure and material intervention (Hacking, 1983).

\section{Hybrid Validation Architectures}

The critique advanced here does not entail abandoning consequence-based evaluation. Transferability across related problems remains a valuable heuristic for identifying promising candidates.

A more robust framework embeds consequence-based scoring within a hierarchy of constraint-sensitive tests. Generativity operates as an exploratory mechanism, while constraint-indexed validation determines whether the resulting structures remain stable across scales, interventions, and embodiments.

Such layered validation resembles practices already present in science, where rapid theoretical exploration is followed by slower empirical consolidation.

\section{Geometry and Constrained Learning Dynamics}

These considerations suggest a reinterpretation of learning itself. Standard optimization treats parameter space as an undifferentiated Euclidean domain. A constraint-aware perspective instead models learning as motion restricted to an admissible manifold.

Information geometry provides a mathematical language for such constrained dynamics, describing updates projected onto structures that preserve relevant invariants (Amari, 2016). Learning becomes alignment with structure rather than unconstrained exploration.

Under this view, elaborate evaluation schemes compensate for an overly permissive search space. If admissible trajectories were constrained from the outset, many formally generative but meaningless candidates would never arise.

\section{Conclusion}

Oracle-free evaluation demonstrates that useful signals can be extracted even when ground truth is unavailable. Its success highlights the power of consequence as an organizing principle within closed formal systems.

Yet consequence alone cannot bear the full weight of validation in open, embodied, and temporally evolving domains. Correctness in such settings emerges from sustained compatibility with constraint rather than from generativity alone.

A comprehensive theory of evaluation must therefore integrate both dimensions. Generativity provides exploration. Constraint provides anchoring. Scientific knowledge advances only where these forces remain in tension, and where explanations are compelled to withstand not only internal coherence but the resistant structure of the world.

\newpage
\begin{thebibliography}{99}

\bibitem{son2026}
J. Son, A. Radhakrishnan, M. S. Alberti, et al.
\newblock Consequence-Based Utility: Oracle-Free Evaluation of Mathematical Reasoning.
\newblock arXiv preprint arXiv:2602.06291, 2026.

\bibitem{kuhn1962}
T. S. Kuhn.
\newblock \emph{The Structure of Scientific Revolutions}.
\newblock University of Chicago Press, 1962.

\bibitem{lakatos1976}
I. Lakatos.
\newblock \emph{Proofs and Refutations}.
\newblock Cambridge University Press, 1976.

\bibitem{merton1973}
R. K. Merton.
\newblock \emph{The Sociology of Science}.
\newblock University of Chicago Press, 1973.

\bibitem{shapin1994}
S. Shapin.
\newblock \emph{A Social History of Truth}.
\newblock University of Chicago Press, 1994.

\bibitem{cartwright1983}
N. Cartwright.
\newblock \emph{How the Laws of Physics Lie}.
\newblock Oxford University Press, 1983.

\bibitem{anderson1972}
P. W. Anderson.
\newblock More Is Different.
\newblock \emph{Science}, 177(4047):393--396, 1972.

\bibitem{hacking1983}
I. Hacking.
\newblock \emph{Representing and Intervening}.
\newblock Cambridge University Press, 1983.

\bibitem{pearl2009}
J. Pearl.
\newblock \emph{Causality}.
\newblock Cambridge University Press, 2009.

\bibitem{polanyi1966}
M. Polanyi.
\newblock \emph{The Tacit Dimension}.
\newblock Doubleday, 1966.

\bibitem{sennett2008}
R. Sennett.
\newblock \emph{The Craftsman}.
\newblock Yale University Press, 2008.

\bibitem{crawford2009}
M. B. Crawford.
\newblock \emph{Shop Class as Soulcraft}.
\newblock Penguin, 2009.

\bibitem{cover2006}
T. Cover and J. Thomas.
\newblock \emph{Elements of Information Theory}.
\newblock Wiley, 2006.

\bibitem{lessig1999}
L. Lessig.
\newblock \emph{Code and Other Laws of Cyberspace}.
\newblock Basic Books, 1999.

\bibitem{marres2017}
N. Marres.
\newblock \emph{Digital Sociology}.
\newblock Polity Press, 2017.

\bibitem{amari2016}
S. Amari.
\newblock \emph{Information Geometry and Its Applications}.
\newblock Springer, 2016.

\bibitem{goodhart1984}
C. A. E. Goodhart.
\newblock Problems of Monetary Management.
\newblock Macmillan, 1984.

\end{thebibliography}

\end{document}
